\scanpage{395}\mspage{309}

\deleted{\unsure{d|c}}\unsure{c'|d'}) Comme d), avec $I$ ens. ordonné.

\uline{Corollaire} 4.13.16. Les hypothèses \inserted{sur $M,\pi_0,\underline S$} étant celles de 4.13.15, pour qu'il existe des cardinaux $\pi>\pi'_0\geq\pi_0$ satisfaisant les conditions $\pi^{\pi'_0}=\pi$ \inserted{(a fortiori $\pi^{\pi_0}=\pi$)} et les conditions équivalentes a) à d') de 4.13.15, il (f et) s que $\underline S$ soit $\deleted{\pi}$accessible dans $M$. De façon plus précise, \deleted{\unsure{En}} si $\underline S$ est accessible, dans $M$, soit $\pi'_0\geq\pi_0$ un cardinal tel que $\underline S$ soit $\pi'_0$-accessible dans $M$, et soit $\pi$ un cardinal \inserted{$>\pi'_0$} tel que $\pi^{\pi'_0}=\pi$. Alors $(\underline S,\pi'_0,\pi)$ satisfait aux conditions équivalentes de 4.13.15.

En effet, le ``il faut'' est évident sur la condition a) de 4.13.15. Pour le ``il suffit'' on \unsure{note|voit} que $\underline S$ $\pi'_0$-accessible dans $M$ implique que $\underline S$ est stable dans $M$ par $\underset{\longrightarrow}{\lim}$ \unsure{grdes|grde} devant $\pi'_0$. Donc pour vérifier a), il suffit de vérifier que $\underline S$ est $\pi$-accessible dans $M$, Or comme $\underline S$ est

% archive sheet number: 394
